\subsection{Kirby calculus for closed oriented four-manifolds}
\label{subsec:topological-input}

This subsection reviews textbook-level handle and Kirby calculus.  We follow
\cite[Chapters~4 and~5]{GompfStipsicz1999}; the original framed-link calculus
is due to Kirby \cite{Kirby1978}.  Throughout, four-manifolds are closed,
connected, oriented, and smooth.  We do not impose a unique minimum or
maximum on a handle decomposition.

\begin{definition}[Handle attachment]
\label{def:handle-attachment}
Let \(W\) be an oriented smooth four-manifold with boundary.  An
\emph{index-\(k\) handle attachment} is a smooth embedding
\[
  \phi\colon S^{k-1}\times D^{4-k}\hookrightarrow\partial W.
\]
It produces \(W\cup_{\phi}(D^k\times D^{4-k})\).  Several handles of the
same index are attached along a finite family of embeddings with pairwise
disjoint images.
\end{definition}

\begin{definition}[Kirby datum]
\label{def:kirby-datum}
A \emph{Kirby datum} \(\mathcal H\) consists of finite handle sets
\(I_k(\mathcal H)\), \(0\leq k\leq4\), and the following attaching data.
Begin with
\[
  W^{(0)}:=\bigsqcup_{q\in I_0(\mathcal H)}B_q^4.
\]
For \(k=1,2,3\), attach the handles successively along pairwise disjoint
embeddings
\[
  \phi_h^k\colon S^{k-1}\times D^{4-k}
  \hookrightarrow\partial W^{(k-1)},
  \qquad h\in I_k(\mathcal H),
\]
and denote the resulting manifold by \(W^{(k)}\).  Finally, for every
\(q\in I_4(\mathcal H)\), choose an orientation-reversing diffeomorphism
\[
  \gamma_q\colon \partial B_q^4\xrightarrow{\ \cong\ }Y_q,
\]
where the \(Y_q\) are precisely the connected components of
\(\partial W^{(3)}\).  Thus the closing maps exhaust the boundary, with one
map for each boundary component.  The resulting closed manifold is required
to be connected.
\end{definition}

\begin{definition}[Realization and presentation]
\label{def:kirby-presentation}
The \emph{realization} of \(\mathcal H\) is
\[
  |\mathcal H|:=W^{(3)}\cup_{\bigsqcup_q\gamma_q}
    \bigsqcup_{q\in I_4(\mathcal H)}B_q^4.
\]
A \emph{Kirby presentation} of a closed connected oriented smooth
four-manifold \(X\) is a pair
\[
  (\mathcal H,\rho),\qquad
  \rho\colon|\mathcal H|\xrightarrow{\ \cong\ }X,
\]
with \(\rho\) orientation preserving.  The choice of any individual closing
map does not affect the oriented diffeomorphism type: two such maps differ by
an orientation-preserving diffeomorphism of \(S^3\), which extends over
\(B^4\) by Cerf's theorem \(\Gamma_4=0\) \cite{Cerf1968}.
\end{definition}

\begin{definition}[Derived handle chain complex]
\label{def:derived-handle-chain-complex}
Put \(W^{(-1)}=\varnothing\) and \(W^{(4)}=|\mathcal H|\), and define
\[
  C_k(\mathcal H):=
  H_k\bigl(W^{(k)},W^{(k-1)};\mathbb Z\bigr).
\]
After orienting the handle cores, their classes identify
\[
  C_k(\mathcal H)\cong
  \bigoplus_{h\in I_k(\mathcal H)}\mathbb Z[h].
\]
The boundary maps are the connecting maps of the handle filtration; their
matrix entries are the signed intersections of attaching and belt spheres.
In particular,
\[
  C_3(\mathcal H)\xrightarrow{\partial_3}
  C_2(\mathcal H)\xrightarrow{\partial_2}
  C_1(\mathcal H).
  \tag{3.1}\label{eq:handle-complex}
\]
Write \(h_k(\mathcal H)=|I_k(\mathcal H)|\).  The chain complex and the
numbers \(h_k\) are derived from the handle attachments; they are not extra
parts of a Kirby datum.
\end{definition}

\begin{definition}[One-ended Kirby datum]
\label{def:one-ended-kirby-datum}
A Kirby datum is \emph{one-ended} if
\(h_0(\mathcal H)=h_4(\mathcal H)=1\).  For such a datum, the \(1\)- and
\(2\)-handles admit the usual reduced record as a finite dotted--framed link
\[
  D_{\mathcal H}=D_{\mathrm{dot}}\sqcup D_{\mathrm{fr}}\subset S^3.
\]
Each dotted component records a \(1\)-handle and each framed component records
a \(2\)-handle.  This link is a convenient record of the full attaching
data, not the definition of a general Kirby datum.
\end{definition}

\begin{proposition}[Existence and one-ended reduction]
\label{prop:kirby-presentation-existence}
Every closed connected oriented smooth four-manifold admits a Kirby
presentation.  Moreover, every Kirby presentation can be changed, by
same-index slides and cancellations of \(0\)--\(1\) and \(3\)--\(4\) pairs,
to a one-ended Kirby presentation of the same manifold.
\end{proposition}

\begin{proof}
A Morse function gives the first assertion.  Since the realized manifold is
connected, the incidence graph of its \(0\)- and \(1\)-handles is connected.
Choose a spanning tree, slide the corresponding \(1\)-handles into cancelling
position, and cancel all but one \(0\)-handle.  Apply the same argument to the
turned-over decomposition to cancel all but one \(4\)-handle.  This is the
standard reduction to a handle decomposition with one minimum and one maximum;
see \cite[Sections~4.2 and~5.1]{GompfStipsicz1999}.
\end{proof}

\begin{definition}[Elementary Kirby--Cerf move]
\label{def:kirby-move}
An elementary Kirby--Cerf move \(m\colon\mathcal H\to\mathcal H'\) is a
level isotopy, a same-index handle slide, or the birth or death of a
cancelling adjacent-index pair.  In dimension four these are the following
thirteen directed types:
\begin{enumerate}[label=(\roman*),ref=\roman*]
  \item\label{move:framed-isotopy} level isotopy of attaching data;
  \item\label{move:one-over-one} a \(1/1\) slide;
  \item\label{move:two-over-one} a \(2/1\) slide;
  \item\label{move:two-over-two} a \(2/2\) slide;
  \item\label{move:three-over-three} a \(3/3\) slide;
  \item a \(0\)--\(1\) birth;
  \item a \(0\)--\(1\) death;
  \item\label{move:one-two-birth} a \(1\)--\(2\) birth;
  \item\label{move:one-two-death} a \(1\)--\(2\) death;
  \item\label{move:two-three-birth} a \(2\)--\(3\) birth;
  \item\label{move:two-three-death} a \(2\)--\(3\) death;
  \item a \(3\)--\(4\) birth; or
  \item a \(3\)--\(4\) death.
\end{enumerate}
Here an \(r/s\) slide moves an index-\(r\) handle over a stationary
index-\(s\) handle.  Outside the supporting neighborhood the handle datum is
unchanged.
\end{definition}

\begin{convention}[Standard local models]
\label{conv:referenced-kirby-dictionary}
The moves in \cref{def:kirby-move} have their standard handle-calculus
meanings.  In a one-ended record, the \(1/1\), \(2/1\), and \(2/2\) moves are
respectively the dotted-circle slide, passage through a dotted handle, and
oriented framed band sum of \cite[Figures~5.39 and~5.36 and
Section~5.1]{GompfStipsicz1999}.  Births insert isolated standard cancelling
pairs and deaths are their inverses.  This convention fixes only the topology
of source and target.
\end{convention}

\begin{theorem}[Kirby--Cerf completeness]
\label{prop:handle-movie-completeness}
Let \((\mathcal H_0,\rho_0)\) and \((\mathcal H_1,\rho_1)\) be Kirby
presentations of \(X_0\) and \(X_1\), and let
\(f\colon X_0\to X_1\) be an orientation-preserving diffeomorphism.  After
pulling the second presentation back along \(f\), the two handle
decompositions of \(X_0\) are related by a finite movie of the thirteen moves
in \cref{def:kirby-move}.  The map \(f\) is the final endpoint
identification; no claim is made that the movie alone realizes its mapping
class.
\end{theorem}

\begin{proof}
This is the standard Cerf-theoretic completeness theorem for handle
decompositions: a generic one-parameter family of Morse functions produces
level isotopies, same-index slides, and births or deaths of adjacent-index
cancelling pairs.  See \cite{Cerf1970} and
\cite[Theorem~4.2.12]{GompfStipsicz1999}.  Pulling back along \(f\) places the
two decompositions on the same manifold, to which that theorem applies.
The endpoint is then identified with the original target by \(f\).
\end{proof}

\begin{remark}[One-ended movies and mapping classes]
\label{rem:mapping-class-refinement}
If both endpoints are one-ended, connected Cerf theory permits a movie that
avoids \(0\)--\(1\) and \(3\)--\(4\) births
\cite[Remark~2.7 and Theorem~3.4]{GayWehrheimWoodward2012}.  This is an
object-level statement.  It does not assert that a based movie realizes an
arbitrarily prescribed element of \(\pi_0\operatorname{Diff}^+(X)\), nor a
complete two-parameter calculus of movies of movies.  Related partial results
for \(S^4\) appear in \cite{Gay2025}; no such refinement is used here.
\end{remark}
